% !TeX program = lualatex
% Source-PDF: 比赛 CONTEST/2021“MINIEYE杯”中国大学生算法设计超级联赛/2021“MINIEYE杯”中国大学生算法设计超级联赛（10）/2021“MINIEYE杯”中国大学生算法设计超级联赛（10）-题解.pdf
\documentclass[11pt]{article}
\usepackage{oipl-vn}

\title{Lời giải trận thứ mười HDU Multi-University 2021}
\author{}
\date{}

\begin{document}
\maketitle

\origtitle{Super League of Chinese College Students Algorithm Design 2021, Contest 10}
\sourcepath{contest/hdu-2021-minieye-10-solutions.pdf}

\section{A. Pty loves sequence}

Làm thế nào để đếm các dãy tốt không trùng và không sót?

Giả sử giá trị lớn nhất là $k$. Ta lấy dãy con đầu tiên dạng
$[1,2,\ldots,k-1,k]$. Sau đó xét việc chèn thêm số vào các khoảng trống giữa
các phần tử này. Không khó thấy rằng để dãy con vừa chọn là dãy con xuất
hiện sớm nhất, điều kiện hạn chế duy nhất là: với khoảng trống đứng trước
$x-1$ và $x$ ($1\le x\le k$), trong khoảng đó không được có $x$. Riêng khoảng
trống cuối cùng thì không bị hạn chế.

Vì vậy chỉ cần liệt kê số lượng $i$ của các số được chèn vào khoảng trống
cuối cùng. Theo phương pháp vách ngăn, số cách điền là
\[
  k^i (k-1)^{n-k-i}\binom{n-k-i+k-1}{k-1}.
\]
Đáp án cho câu hỏi thứ nhất là
\[
  \mathrm{Ans}
  =
  \sum_{k=1}^{n}\sum_{i=0}^{n-k}
  k^i (k-1)^{n-k-i}\binom{n-i-1}{k-1}.
\]

Cách làm cho câu hỏi thứ hai có thể suy ra từ câu hỏi thứ nhất. Có một kết
luận khá hiển nhiên:
\begin{quote}
Khi $k$ cố định, số lần xuất hiện của các giá trị $1,\ldots,k$ trong các dãy
tốt là như nhau.
\end{quote}

Chứng minh ngắn gọn như sau. Trong cách làm của câu hỏi thứ nhất, khi $k$ cố
định, hạn chế là khoảng trống thứ $i$ không được điền $i$, còn khoảng trống
cuối cùng tùy ý. Vì ràng buộc trên mỗi phần tử trong $1,\ldots,k$ là giống
nhau, nên bản chất của các phần tử này là như nhau, do đó số lần xuất hiện
cũng phải bằng nhau.

Vì thế ta liệt kê $k$ và cộng vào \texttt{ans[1..k]}:
\[
  \left(
  \sum_{i=0}^{n-k}
  k^i (k-1)^{n-k-i}\binom{n-i-1}{k-1}
  \right)\cdot \frac{n}{k}.
\]

Vấn đề cuối cùng là mô-đun của bài này không phải số nguyên tố, nên không có
nghịch đảo modulo; do đó cần đưa thừa số $n/k$ vào trong biểu thức. Quan sát
công thức, ta nhận được:
\begin{itemize}
  \item nếu $i>0$, chỉ cần đổi $k^i$ thành $k^{i-1}$;
  \item nếu $i=0$, đổi
  $\binom{n-1}{k-1}\cdot \frac{n}{k}$ thành $\binom{n}{k}$.
\end{itemize}

Như vậy trong công thức chỉ còn dùng lũy thừa và tổ hợp. Tổ hợp có thể tiền
xử lý bằng tam giác Pascal, nên mô-đun không cần là số nguyên tố.

Độ phức tạp thời gian là $O(n^2)$.

Thực ra khi $p=998244353$, bài này có thể làm trong
$O(n\sqrt n\log n)$ bằng đa thức và tính giá trị tại nhiều điểm. Vì cách này
có hằng số khá lớn và mã dài, không phù hợp với thi ACM, nên người ra đề đã
không dùng phiên bản đó.

\section{B. Pty with card}

Ta quy ước một vòng rút bài là khi tất cả người còn trong vòng lần lượt rút
bài một lần. Một trạng thái là dãy số lượng bài trên tay của tất cả người
trong vòng, theo thứ tự.

Xét một trạng thái sẽ thay đổi như thế nào. Trạng thái ban đầu là
$(1,1,1,\ldots,1,1)$. Sau một vòng rút bài, trạng thái trở thành
$(3,2,2,2,\ldots,2,2)$ hoặc $(4,2,2,2,\ldots,2)$, và số người còn lại chắc
chắn là
\[
  \left\lfloor \frac{N-1}{2}\right\rfloor .
\]

Nếu một trạng thái còn $2m$ người, trước khi có ai đó bị loại, số bài trên
tay của mỗi người ở mỗi vòng chắc chắn luôn tăng $1$ hoặc luôn giảm $1$.

Nếu một trạng thái còn $2m+1$ người, sau hai vòng, nếu không có ai bị loại,
số bài trên tay của mỗi người đều không đổi. Đây là nguyên nhân tạo ra chu
kỳ; độ dài chu kỳ là $4m+2$.

Thử tay vài dữ liệu, ta thấy vài tính chất:
\begin{enumerate}
  \item Nếu một trạng thái còn $2m$ người, thì sau vòng có người bị loại, số
  người còn lại chắc chắn là $m$. Hơn nữa trạng thái có dạng
  \[
    (P+1,P,P,P,\ldots,P,P)
  \]
  hoặc
  \[
    (P+2,P,P,P,\ldots,P,P),
  \]
  và việc là dạng trước hay sau chỉ phụ thuộc vào tính chẵn lẻ của $N$ ban
  đầu, trong đó $P=2^k,\ k\in\mathbb{Z}^+$.
  \item Sau vòng đầu tiên, nếu một trạng thái còn $2m+1$ người thì chắc chắn
  xuất hiện chu kỳ.
\end{enumerate}

Các tính chất này chứng minh bằng quy nạp khá thuận tiện. Gọi trạng thái ban
đầu của $m$ người là trạng thái còn $m$ người sau vòng sớm nhất. Nếu trạng
thái có dạng $(P+1,P,P,P,\ldots,P,P)$ thì gọi là dạng A; nếu có dạng
$(P+2,P,P,P,\ldots,P,P)$ thì gọi là dạng B, với
$P=2^k,\ k\in\mathbb{Z}^+$.

Bắt đầu xét từ vòng thứ hai. Ta chứng minh mệnh đề $p$ sau: nếu $N$ lẻ thì
mọi trạng thái ban đầu có thể đạt được đều có dạng
$(P+1,P,P,P,\ldots,P,P)$ và lần rút đầu tiên chắc chắn rút $2$ lá; nếu $N$
chẵn thì trạng thái ban đầu có dạng $(P+2,P,P,P,\ldots,P,P)$ và lần rút đầu
tiên chắc chắn rút $1$ lá.

Ngay từ đầu, trạng thái có
$n=\left\lfloor (N-1)/2\right\rfloor$ người và thỏa mệnh đề $p$.

Khi trạng thái có $2m$ người và thỏa mệnh đề $p$, hiển nhiên $P\ge 2$.
\begin{enumerate}
  \item Nếu $N$ lẻ, trạng thái ban đầu là
  $(P+1,P,P,P,\ldots,P,P)$, với $P+1\ge 3$, và lần rút đầu tiên rút $2$ lá.
  Sau mỗi vòng, giá trị ở vị trí lẻ giảm $1$, còn giá trị ở vị trí chẵn tăng
  $1$. Cho đến khi trạng thái trở thành
  \[
    (2,2P-1,1,2P-1,1,2P-1,\ldots,1,2P-1),
  \]
  không có người nào có số bài về $0$. Thêm một vòng nữa, các giá trị ở vị
  trí lẻ đều trở thành $0$ và những người này ra khỏi vòng; trạng thái trở
  thành
  \[
    (2P+1,2P,2P,2P,\ldots,2P,2P).
  \]
  Tiếp theo lại đến lượt rút $2$ lá, nên vẫn thỏa mệnh đề $p$.

  \item Nếu $N$ chẵn, trạng thái ban đầu là
  $(P+2,P,P,P,\ldots,P,P)$, với $P+2\ge 4$, và lần rút đầu tiên rút $1$ lá.
  Sau mỗi vòng, giá trị ở vị trí chẵn giảm $1$, còn giá trị ở vị trí lẻ tăng
  $1$. Cho đến khi trạng thái trở thành
  \[
    (2P+1,1,2P-1,1,2P-1,\ldots,2P-1,1),
  \]
  không có người nào có số bài về $0$. Thêm một vòng nữa, các giá trị ở vị
  trí chẵn đều trở thành $0$ và những người này ra khỏi vòng; trạng thái trở
  thành
  \[
    (2P+2,2P,2P,2P,\ldots,2P,2P).
  \]
  Tiếp theo lại đến lượt rút $1$ lá, nên vẫn thỏa mệnh đề $p$.
\end{enumerate}

Khi trạng thái có $2m+1$ người và thỏa mệnh đề $p$, hiển nhiên $P\ge 2$.
\begin{enumerate}
  \item Nếu $N$ lẻ, trạng thái là $(P+1,P,P,P,\ldots,P,P)$ và trước hết rút
  $2$ lá. Sau một vòng, trạng thái trở thành
  \[
    (P+1,P+1,P-1,P+1,P-1,\ldots,P+1,P-1).
  \]
  Vì $P\ge 2$ nên không có số nào về $0$. Vòng kế tiếp sẽ bắt đầu bằng rút
  $1$ lá, và sau vòng đó trạng thái trở lại
  $(P+1,P,P,P,\ldots,P,P)$, tức là xuất hiện chu kỳ.

  \item Nếu $N$ chẵn, trạng thái là $(P+2,P,P,P,\ldots,P,P)$ và trước hết
  rút $1$ lá. Sau một vòng, trạng thái trở thành
  \[
    (P+2,P-1,P+1,P-1,P+1,\ldots,P-1,P+1).
  \]
  Vì $P\ge 2$ nên không có số nào về $0$. Vòng kế tiếp sẽ bắt đầu bằng rút
  $2$ lá, và sau vòng đó trạng thái trở lại
  $(P+2,P,P,P,\ldots,P,P)$, tức là xuất hiện chu kỳ.
\end{enumerate}

Chứng minh hoàn tất.

Vậy với $N$, nếu $N\le 2$ thì $F(N)=0$. Ngược lại đặt
\[
  m=\left\lfloor\frac{N-1}{2}\right\rfloor,\qquad
  S=\min\left\{\frac{m}{2^k}\mid 2^k\mid m,\ k\in\mathbb{Z}\right\}.
\]
Nếu $S=1$ thì $F(N)=0$, ngược lại $F(N)=2S$.

Sau khi biết điều này, ta xét cách làm bài toán. Vì cần chia cho lũy thừa
lớn nhất của $2$, có thể dùng \texttt{01-trie} để duy trì. Trong phân rã tâm,
với mỗi điểm xuất phát, đưa $v_i$ cộng với khoảng cách đến tâm phân rã vào
\texttt{01-trie}. Với một điểm có khoảng cách đến tâm phân rã là $l$, ta
thực hiện $l$ lần tăng $1$ trên các giá trị trong \texttt{01-trie}, rồi truy
vấn đáp án.

Với thao tác tăng $1$ trên \texttt{01-trie}: điều kiện là các số trong
\texttt{01-trie} được chèn từ bit thấp lên bit cao. Khi cộng $1$ vào một số,
biểu diễn nhị phân thay đổi dạng
\[
  111110\ldots \to 000001\ldots
\]
với bit thấp ở trước; điều này tương đương với hoán đổi các nút $0/1$ trên
\texttt{01-trie}.

Với truy vấn, trước hết giảm $v_i$ đọc vào đi $1$, rồi cộng đáp án ở nhánh
$0$ và nhánh $1$. Một nút trên \texttt{01-trie} đại diện cho tất cả số có
$i$ bit đầu giống nhau; ta cần duy trì tổng của các số đó sau khi chia lấy
sàn cho $2^i$. Khi tính đáp án cần chú ý loại bỏ các lũy thừa của $2$, vì
đáp án của chúng là $0$.

Độ phức tạp thời gian là $O(n\log^2 n)$.

\section{C. Pty loves lines}

\subsection*{Cách 1}

Ý nghĩa bài toán là tìm mọi giá trị có thể có của
\[
  \frac{n(n-1)}{2}
  -
  \sum_i \frac{a_i(a_i-1)}{2},
\]
với điều kiện $\sum_i a_i=n$.

Vế trái có phần hằng số, nên ta xét mọi giá trị có thể đạt được của vế phải.
Vì $a_i=1$ có đóng góp $0$, nếu một phương án với $\sum_i a_i=x$ có thể tạo
ra giá trị $s$, thì tổng lớn hơn $x$ cũng có thể tạo ra giá trị đó.

Chỉ cần đặt $f_n$ là giá trị nhỏ nhất của $\sum_i a_i$ để vế phải bằng $n$.
Chuyển trạng thái tương tự quy hoạch động ba lô. Độ phức tạp thời gian là
$O(n^3)$.

\subsection*{Cách 2}

Đặt $f[i][j]$ biểu thị việc dùng $i$ đường thẳng có thể tạo ra $j$ giao điểm
hay không.

Có thể dùng \texttt{bitset} để tối ưu chuyển trạng thái, đạt
$O(n^4/w)$, nhưng vẫn khó qua giới hạn thời gian $1$ giây.

Quan sát thấy đáp án có một hậu tố khả thi liên tiếp khá dài. Sau khi lập
bảng, ta thấy số giao điểm không khả thi lớn nhất không vượt quá $35000$, vì
vậy thời gian có thể tối ưu thành $35000\times n^2/w$ và đủ qua bài này.

\section{D. Pty hates prime numbers}

\subsection*{Cách 1}

Xét bao hàm--loại trừ kinh điển. Liệt kê mỗi số được chọn hay không được
chọn; giả sử tích các số được chọn là $S$, khi đó
\[
  \mathrm{ans} \mathrel{+}= \left\lfloor\frac{n}{S}\right\rfloor
  \cdot (-1)^{\text{số lượng số được chọn}}.
\]
Độ phức tạp thời gian là $O(2^k)$.

\subsection*{Cách 2}

Giả sử $k=8$. Tích của $8$ số nguyên tố đầu tiên là
\[
  2\cdot 3\cdot 5\cdot 7\cdot 11\cdot 13\cdot 17\cdot 19
  =9699690,
\]
rất nhỏ.

Chú ý rằng
\[
  [x\text{ không là bội của }8\text{ số nguyên tố đầu}]
  =
  [x\bmod 9699690\text{ không là bội của }8\text{ số nguyên tố đầu}].
\]
Do đó
\[
  \text{đáp án của }n
  =
  (\text{đáp án của }9699690)
  \left\lfloor\frac{n}{9699690}\right\rfloor
  +(\text{đáp án của }n\bmod 9699690).
\]
Tiền xử lý đáp án cho $n=0\sim 9699690$ là có thể trả lời mỗi truy vấn trong
$O(1)$.

\subsection*{Cách 3}

Kết hợp hai cách trên:
\begin{enumerate}
  \item Khi $k\le 8$, làm tùy ý bằng cách trên.
  \item Khi $k>8$, trước hết bao hàm--loại trừ từ số nguyên tố thứ $9$ đến
  số nguyên tố thứ $k$; khi còn lại $8$ số nguyên tố đầu thì dùng trực tiếp
  đáp án đã tiền xử lý trong $O(1)$.
\end{enumerate}

Độ phức tạp thời gian là
\[
  O\left(\sum_{i=1}^{8}\frac{9699690}{p[i]}+T\cdot 2^8\right),
\]
dễ dàng qua bài này.

\section{E. Pty loves book}

Đưa tất cả chuỗi vào cùng một automaton AC. Xét việc tính, với mỗi nút $x$
của automaton AC đại diện cho chuỗi $s$,
\[
  \sum_{i=1}^{|s|} f(s,i,|s|)^5.
\]

Đặt $\mathrm{tot}[x][k]$ là tổng lũy thừa bậc $k$ của giá trị của tất cả hậu
tố của chuỗi mà $x$ đại diện. Đặt $\mathrm{failtot}[x][k]$ là tổng lũy thừa
bậc $k$ của các hậu tố của chuỗi mà $x$ đại diện, nhưng có đầu trái nằm trước
\(\operatorname{fail}(x)\).

$\mathrm{tot}[x]$ được chuyển từ $\mathrm{tot}$ của
\(\operatorname{fail}(x)\) và $\mathrm{failtot}[x]$.

Xét cách tính $\mathrm{failtot}[x]$. Ta thấy đây chính là tổng
$\mathrm{failtot}$ của tất cả nút đi qua khi $x$ tìm liên kết fail, bắt đầu
từ cha của $x$ và bao gồm cả cha của $x$, cộng với phần giá trị tăng thêm do
thêm một ký tự, tức là tổng giá trị của tất cả fail của $x$. Cần đặc biệt
tính riêng giá trị của toàn bộ chuỗi tại $x$.

Để xử lý lũy thừa bậc $5$, duy trì thêm tổng lũy thừa bậc $0$ đến $4$ là có
thể gộp nhanh.

Độ phức tạp thời gian là
\[
  O\left(\left(|S|+\sum |T|\right)\times 5^2\right).
\]

\section{F. Pty loves lcm}

Khi $y-x+1>2$,
\[
  f(x,y)\approx \frac{x^{y-x+1}}{(y-x)!},
\]
nên số lượng cặp có $f(x,y)<10^{18}$ xấp xỉ $10^6$; có thể tiền xử lý bạo
lực.

Với $y-x=1$, ta có
\[
  \gcd(x,y)=\gcd(x,x+1)=1,
\]
nên bài toán trở thành tính các tổng dạng
\[
  \sum_{i=1}^{n}\varphi(i)\varphi(i+1).
\]
Vì $n\le 10^9$, có thể lập bảng theo đoạn.

Cụ thể, cứ mỗi $10^6$ giá trị lập một bảng, khi đó kích thước bảng khoảng
$10\,\mathrm{kb}$.

Cả cách lập bảng theo đoạn lẫn lời giải chính đều dùng sàng theo đoạn. Nếu
sàng đoạn $[l,r]$, trước hết xử lý các số nguyên tố không vượt quá
$\sqrt r$; với mỗi số nguyên tố, liệt kê các bội của nó trong $[l,r]$, sàng
đi và nhân với giá trị $\varphi$ tương ứng. Cuối cùng nhân tiếp với giá trị
$\varphi$ của phần nguyên tố còn lại sau khi sàng.

Độ phức tạp thời gian là $O((r-l)\log\log r)$.

\section{G. Pty loves graph}

Bài toán tương đương với việc trên một chu trình Hamilton có một số dây
cung. Ta cần bố trí mỗi dây cung ở bên trong hoặc bên ngoài vòng tròn, đồng
thời hai dây cung giao nhau thực sự, tức là không tính trường hợp chứa nhau
hoặc rời nhau, không được cùng nằm bên trong hoặc cùng nằm bên ngoài.

Cắt chu trình Hamilton tại một vị trí bất kỳ, quan hệ giao nhau thực sự vẫn
được giữ nguyên. Xem mỗi dây cung là một đỉnh và nối cạnh giữa hai đỉnh có
ràng buộc.

Khi đó bài toán tương đương với kiểm tra đồ thị mới có tô được hai màu hay
không. Ta chỉ quan tâm đến tính liên thông và khả năng tô màu của đồ thị.
Vì vậy, sắp theo đầu trái giảm dần và đầu phải tăng dần, rồi liệt kê đoạn
bên phải trong hai đoạn giao nhau. Đồng thời duy trì tập các đầu phải của
những đoạn có đầu trái nhỏ hơn đầu trái của đoạn hiện tại. Ta thấy các cạnh
cần thêm ứng với một đoạn liên tiếp của các đầu phải.

Trong đồ thị cuối cùng, đoạn liên tiếp này là liên thông và cùng màu. Vì thế
khi về sau lại nối vào đoạn này, chỉ cần thêm một cạnh. Hơn nữa các đầu phải
chỉ bị xóa chứ không được chèn thêm, nên có thể dùng DSU và \texttt{set} để
tối ưu quá trình thêm cạnh. Mỗi lần, hợp nhất toàn bộ các đầu phải trong
đoạn cần nối cạnh. Cuối cùng kiểm tra đồ thị có phải đồ thị hai phía hay
không.

Độ phức tạp thời gian là $O(n\log n)$, bộ nhớ $O(n)$. Cách dùng cây đoạn có
hằng số thời gian và bộ nhớ tốt cũng có thể qua.

\section{H. Pty loves string}

Xét một vị trí xuất hiện $S[l..r]$. Nó thỏa
$S[l..l+x-1]$ bằng $S[1..x]$, $S[r-y+1..r]$ bằng $S[n-y+1..n]$, và
\[
  l+x=r-y+1.
\]

Hai điều kiện đầu gợi đến border; dùng KMP là có thể tìm các $x$ hợp lệ cho
mọi $l+x-1$.

Sau khi chạy KMP, nối một cạnh từ $i$ đến border của $i$. Ta thu được một
cây; mọi vị trí xuất hiện của một $x$ chính là cây con của nó.

Vậy bài toán chuyển thành: cho hai cây, hỏi trong hai cây con có bao nhiêu
nút có cùng chỉ số?

Bài toán này có thể chuyển thành đếm điểm hai chiều, rồi duy trì bằng cấu
trúc dữ liệu, chẳng hạn quét đường và cây Fenwick.

Độ phức tạp là $O(n\log n)$.

\section{I. Pty loves SegmentTree}

Gọi $f_n$ là đáp án khi khoảng của nút gốc là $[1,n]$. Trước hết có
$f_1=1$. Sau đó xét nút gốc thuộc loại A hay loại B, ta có
\[
  f_n=(A-B)f_{n-k}f_k+B\sum_{i=1}^{n-1} f_i f_{n-i}.
\]

Ta tính $f_k$ trước, rồi đặt $A$ mới bằng $(A-B)f_k$. Xét hàm sinh $F(x)$
của dãy $\{f_n\}$. Dễ dàng có
\[
  F(x)=Ax^kF(x)+BF^2(x)+x.
\]
Giải phương trình thu được
\[
  F(x)=\frac{1-Ax^k\pm\sqrt{A^2x^{2k}-2Ax^k-4Bx+1}}{2B}.
\]
Vì hằng số tự do của $F(x)$ là $0$, ta chọn dấu trừ.

Ta thử tránh ảnh hưởng của căn bậc hai. Đặt
\[
  Q^2(x)=A^2x^{2k}-2Ax^k-4Bx+1.
\]
Trước hết ta có
\[
  2BF(x)=1-Ax^k-Q(x). \tag{1}
\]
Lấy đạo hàm hai vế:
\[
  2BF'(x)=-kAx^{k-1}-Q'(x),
\]
suy ra
\[
  Q'(x)=-kAx^{k-1}-2BF'(x). \tag{2}
\]

Quay lại định nghĩa,
\[
  Q^2(x)=A^2x^{2k}-2Ax^k-4Bx+1.
\]
Lấy đạo hàm hai vế:
\[
  2Q'(x)Q(x)=2kA^2x^{2k-1}-2kAx^{k-1}-4B,
\]
hay
\[
  Q'(x)Q(x)=kA^2x^{2k-1}-kAx^{k-1}-2B.
\]
Thay giá trị $Q'(x)$ từ (2), ta được
\[
  (-kAx^{k-1}-2BF'(x))Q(x)
  =
  kA^2x^{2k-1}-kAx^{k-1}-2B.
\]
Nhân hai vế với $Q(x)$:
\[
  (Akx^{k-1}+2BF'(x))Q^2(x)
  =
  Q(x)(2B+kAx^{k-1}-kA^2x^{2k-1}).
\]

Từ (1) ta biết $Q(x)$, còn $Q^2(x)$ là định nghĩa. Thay vào:
\[
\begin{aligned}
 &(Akx^{k-1}+2BF'(x))(A^2x^{2k}-2Ax^k-4Bx+1)\\
 &\quad =(1-Ax^k-2BF(x))(2B+kAx^{k-1}-kA^2x^{2k-1}).
\end{aligned}
\]

Khai triển:
\[
\begin{aligned}
&2BF'(x)(A^2x^{2k}-2Ax^k-4Bx+1)
  +Akx^{k-1}+A^3kx^{3k-1}\\
&\quad -2A^2kx^{2k-1}-4AkBx^k\\
&=
-2BF(x)(2B+kAx^{k-1}-kA^2x^{2k-1})
  +2B+kAx^{k-1}\\
&\quad -kA^2x^{2k-1}-2ABx^k-kA^2x^{2k-1}+kA^3x^{3k-1}.
\end{aligned}
\]

Rút gọn:
\[
  F'(x)(A^2x^{2k}-2Ax^k-4Bx+1)-2Akx^k
  =
  -F(x)(2B+kAx^{k-1}-kA^2x^{2k-1})-Ax^k+1.
\]

Lấy hệ số của $x^n$ ở hai vế:
\[
\begin{aligned}
 &(n+1)f_{n+1}
 +A^2(n-2k+1)f_{n-2k+1}
 -2A(n-k+1)f_{n-k+1}
 -4Bnf_n\\
 &\quad =
 -2Bf_n-kAf_{n-k+1}+kA^2f_{n-2k+1}
 +[n=k](A(2k-1)).
\end{aligned}
\]

Sắp xếp lại:
\[
\begin{aligned}
  (n+1)f_{n+1}
  &=2B(2n-1)f_n
    +A(2n-3k+2)f_{n-k+1}\\
  &\quad
    -A^2(n-3k+1)f_{n-2k+1}
    +[n=k](A(2k-1)).
\end{aligned}
\]

Đến đây ta đã có công thức truy hồi của $f_{n+1}$ theo
$f_n$, $f_{n-k+1}$ và $f_{n-2k+1}$, nên chỉ cần truy hồi trực tiếp.

Độ phức tạp là $O(n)$.

\section{J. Pty plays game}

Nếu mỗi bên chỉ có một người, chỉ cần so sánh thời gian hai người đánh chết
đối phương:
\[
  t=\min(h_1/d_2,\ h_2/d_1).
\]
Lính $2$ thắng khi và chỉ khi
\[
  h_1/d_2<h_2/d_1,
\]
chuyển vế được
\[
  h_1d_1<h_2d_2.
\]

Giả sử thời gian hai người đánh nhau là $t$, khi đó
\[
  t=\min(h_1/d_2,\ h_2/d_1).
\]
Giả sử giá trị này là $h_1/d_2$. Sau thời gian $t$, máu của hai người lần
lượt là
\[
  h_1-(h_1/d_2)d_2,\qquad h_2-(h_1/d_2)d_1.
\]
Bây giờ lính $2$ phải đánh với người tiếp theo. Xét tích công kích với máu
của lính này:
\[
  (h_2-(h_1/d_2)d_1)d_2=h_2d_2-h_1d_1.
\]

Nếu xem tích công kích và máu của một lính là ``sức chiến đấu'' của lính đó,
thì sau khi hai lính đánh nhau, sức chiến đấu của người có sức chiến đấu cao
hơn sẽ giảm đi sức chiến đấu của người thấp hơn.

Cuối cùng đội còn lại chắc chắn là đội có tổng sức chiến đấu cao hơn.

Vì công kích và máu đều được viết dưới dạng hàm bậc nhất theo $x$, sức chiến
đấu của một đội là một hàm bậc hai. Bài toán trở thành truy vấn số nguyên
không âm nhỏ nhất $x$ sao cho hàm bậc hai đó dương.

Bài này có thể bị kẹt độ chính xác, nên cần dùng một dạng số chính xác cao
hoặc số nguyên lớn.

Độ phức tạp thời gian là $O(n)$.

\end{document}
